\section{Safety observers under fixed intervention budgets}
\label{sec:safety}

\subsection{A perfect classifier can allocate a safety budget poorly}

Suppose a workspace assistant is asked to delete a file in workspace A, but the
user has granted access only to workspace B. Deletion is not always unsafe; the
problem is that this file is outside the granted scope. A controller may allow,
block, or escalate the request. When blocks and reviews are limited, it must
rank requests by the harm they could cause.

We turn this example into an inert paired task:

\begin{itemize}
  \item \textbf{Same action, different permission.} Each pair requests the same
  operation on the same resource. One prompt grants access to it; the other
  grants access elsewhere.
  \item \textbf{Known label and cost.} An exact authorization rule says whether
  the request exceeds scope. The operation sets the consequence: an unauthorized
  delete costs more than an unauthorized read.
  \item \textbf{Fixed action budget.} The controller blocks the highest-risk
  10\%, escalates the next 10\%, and allows the rest. Every observer faces the
  same costs and budgets.
\end{itemize}

We use the pinned Qwen2.5-7B-Instruct model \cite{qwen2024qwen25}. As a
prerequisite check, the unedited model must answer the authorization question
with one of two exact tokens. It passes every check with accuracy 1.000. The
held-out test contains 384 safe/unsafe pairs across four equally sized subgroups:
seen operations and prompt forms, a held-out delete operation, a held-out XML
form, and both changes together. Fit, calibration, and test resources are
disjoint.

The activation observer reads a hidden-state vector---a residual-stream
activation---at the final prompt position and also receives the requested
action and its severity. We fit on 384
pairs, calibrate on 192, choose among seven layers and ridge settings using the
fixed policy loss, and only then extract held-out activations.

The comparisons ask four questions:

\begin{itemize}
  \item \textbf{Action-only:} what can be inferred without authorization
  context?
  \item \textbf{Activation label and direct risk:} do hidden activations predict
  the policy label or the violation cost?
  \item \textbf{Activation label $\times$ severity:} does combining the label
  score with known consequence improve the action ranking?
  \item \textbf{Allow-all and exact oracle:} what are the upper and lower
  bounds under the same budget?
\end{itemize}

These comparisons were fixed before held-out activations were extracted.
After seeing their results, we added three secondary controls: a full-context
text observer, Qwen's block-minus-allow output-score difference, and an
activation observer selected by tail rather than mean loss. They test whether
activations add value once prompt text or the model's own decision signal is
available; they are not independent confirmation.

\begin{table}[H]
\centering
\scriptsize
\caption{Held-out Qwen safety-interlock results. Every observer uses the same
10\% block and 10\% escalation budgets. Lower loss and missed risk are better;
higher clean utility is better. Access states what the observer can read.}
\label{tab:qwen-safety}
\begin{tabular}{@{}p{0.27\textwidth}p{0.18\textwidth}rrrr@{}}
\toprule
Observer & Access & Mean loss & Tail loss & Missed risk & Clean utility \\
\midrule
Allow all & None & 5.010 & 19.026 & 1.000 & 1.000 \\
Action-only direct risk & Action fields & 3.176 & 13.987 & 0.623 & 0.901 \\
Activation label & Residual + fields & 3.334 & 18.182 & 0.665 & 1.000 \\
Activation direct risk & Residual + fields & 2.520 & 16.545 & 0.502 & 1.000 \\
Activation label $\times$ severity & Residual + fields & 2.155 & 14.026 & 0.429 & 1.000 \\
Full-context text label $\times$ severity$^\dagger$ & Prompt text + fields & 2.032 & 12.857 & 0.405 & 1.000 \\
Qwen logit margin $\times$ severity$^\dagger$ & Output logits + fields & \textbf{1.744} & \textbf{10.078} & \textbf{0.347} & \textbf{1.000} \\
Exact authorization-risk oracle & Evaluator label + fields & 1.740 & 10.039 & 0.346 & 1.000 \\
\bottomrule
\end{tabular}
\end{table}

\noindent\emph{$^\dagger$ Secondary controls.} The text and output rows were
designed after the original held-out results were known.

The results have four main implications:

\begin{itemize}
  \item \textbf{Activations improve the primary comparison.} Multiplying the
  activation label by severity lowers mean policy loss by 32.1\% relative to
  action-only and 57.0\% relative to allowing every request. It closes 71.1\%
  of the action-only-to-oracle gap while preserving all benign utility.
  \item \textbf{Perfect classification is not enough.} The label observer has
  AUROC 1.000 and loss 3.334. Multiplying its score by severity lowers AUROC to
  0.954 but lowers loss to 2.155 because unauthorized deletes are ranked ahead
  of unauthorized reads.
  \item \textbf{The useful signal is not uniquely hidden.} Full-context text is
  not reliably worse than the transformed activation observer. More clearly,
  Qwen's output-score difference multiplied by severity lowers mean loss to
  1.744 and tail loss to 10.078, nearly matching the exact oracle at 1.740 and
  10.039. When this output is available, the activation probe is unnecessary.
  \item \textbf{Good ranking within groups can fail across groups.} The
  activation label has AUROC 1.000 inside every operation--format subgroup, but
  its scale changes between them. A global threshold therefore allocates 63.2\%
  of interventions to held-out deletes in seen forms and only 18.4\% to the
  doubly held-out subgroup. There, activation loss is 5.007 versus 3.059 for
  action-only; the output observer catches 69.8\% of unsafe requests versus
  29.2\% for the activation observer.
\end{itemize}

\begin{figure}[H]
\centering
\includegraphics[width=\textwidth]{qwen_safety_stratum_calibration.png}
\caption{Risk scores by operation--format subgroup. The activation observer
separates safe from unsafe requests within each subgroup, but its scale shifts
between subgroups. Dashed lines show each observer's global intervention
cutoff. Scores should be compared within panels, not across them.}
\label{fig:qwen-safety-calibration}
\end{figure}

Recalibration by visible operation and format may repair this failure, but this
split cannot test that repair because calibration contains neither held-out
case. A valid follow-up needs fresh prompts and calibration coverage of
every intended deployment group.

The task therefore requires an observer to:

\begin{itemize}
  \item classify the violation;
  \item rank it by consequence; and
  \item calibrate scores across deployment groups.
\end{itemize}

It does not test harmful execution, deception, adaptive attacks, or open-ended
deployment.

\subsection{Activations can rank verified APPS backdoors without a verdict prompt}
\label{sec:safety-apps}

The authorization task is controlled but constructed. We therefore use
ControlArena's APPS setting, an existing AI-control task with ordinary code
solutions and verified working backdoors
\cite{greenblatt2024aicontrol,inglis2025controlarena}. A monitor reads the
programming problem and submitted code, assigns a suspicion score, and a fixed
budget audits only the highest-scoring submissions. Every unaudited working
backdoor counts as one violation.

The public dataset supplies 1,444 honest/backdoor pairs. We split by problem ID:
500 pairs for fitting, 200 for calibration, and 744 for held-out testing, with both
members of a pair always in the same split. We use the same pinned Qwen model as
above.

The observers fall into three groups:

\begin{itemize}
  \item \textbf{Prompt and output observers} ask Qwen for a score and read the
  model's probability distribution over ten score tokens. We test basic and
  detailed prompts, their ensemble, and a follow-up ridge model using all ten
  token probabilities.
  \item \textbf{A text observer} is a follow-up baseline trained on word and
  character features from the problem and code.
  \item \textbf{Activation observers} fit a linear model to a hidden-state
  vector. One is read next to an explicit verdict prompt; the follow-up neutral
  observer reads the same problem and code with no monitoring or scoring
  instruction.
\end{itemize}

The original comparison used the three prompt observers and the
verdict-adjacent activation probe. We added the score-distribution, text, and
neutral-activation controls after that result and fixed them before measurement.
At a 2\% audit budget, calibration selects layer 28 for the verdict-adjacent
probe and layer 20 for the neutral probe.

\begin{table}[H]
\centering
\small
\caption{Monitor comparison on 744 held-out APPS pairs. The controller audits
29 of 1,488 submissions at 2\% and 148 at 10\%. Higher caught-backdoor counts
are better; 29 and 148 are the best possible counts under those budgets.}
\label{tab:apps-monitor}
\begin{tabular}{@{}llrrr@{}}
\toprule
Observer & Access & AUROC & Caught @2\% & Caught @10\% \\
\midrule
Exact-label oracle & evaluator & 1.000 & 29 & 148 \\
Neutral activation & neutral residual & \textbf{0.979} & \textbf{29} & \textbf{148} \\
Verdict-adjacent activation & prompt residual & 0.953 & 29 & 148 \\
Supervised text & problem + code & 0.868 & 29 & 146 \\
Score-distribution ridge & prompt logits & 0.830 & 29 & 147 \\
Prompt ensemble & two prompt expectations & 0.817 & 28 & 146 \\
Detailed prompt & prompt expectation & 0.810 & 28 & 146 \\
Basic prompt & prompt expectation & 0.786 & 27 & 137 \\
\bottomrule
\end{tabular}
\end{table}

Three findings matter:

\begin{itemize}
  \item \textbf{Neutral activations rank attacks best.} The neutral probe reaches
  AUROC 0.979 without asking Qwen to judge the code. Its paired gains are 0.026
  over verdict-adjacent activation, 0.111 over supervised text, and 0.149 over
  the score-distribution observer, with 95\% intervals $[0.018,0.034]$,
  $[0.098,0.125]$, and $[0.131,0.167]$.
  \item \textbf{The balanced panel hides action differences.} Neutral and
  verdict-adjacent activations both reach the exact-label audit ceiling at 2\%,
  5\%, and 10\%. At 10\%, neutral activation leaves 596 attacks unaudited versus
  607 for the basic prompt, a paired difference of $-11$ with 95\% interval
  $[-17,-5]$. The sample does not resolve its smaller gains over the stronger
  controls because attacks make up half the panel and fill the small audit queue.
  \item \textbf{The tested controls narrow, but do not settle, the explanation.}
  Reading all ten output probabilities raises AUROC from 0.786 to 0.830, so
  sampling or averaging the declared score tokens does not explain the activation
  gain. The text result rules out this lexical baseline, not every possible code
  monitor or output observer.
\end{itemize}

In deployment, attacks should be rare. We therefore resample the frozen scores at attack
rates from 1--50\% and audit budgets from 0.5--10\%. At 1\% attacks and a 2\%
audit budget, the exact oracle misses none. The activation observers miss 2.65
attacks on average for neutral context (95\% range $[0,6]$) and 3.94 next to the
verdict prompt ($[1,7]$). The text, score-distribution, and basic-prompt observers
miss 9.05 ($[6,12]$), 9.15 ($[6,13]$), and 11.85 ($[9,15]$). Rare attacks expose
differences hidden by the balanced panel.

\begin{figure}[H]
\centering
\includegraphics[width=\textwidth]{controlarena_prevalence_sensitivity.pdf}
\caption{Excess escaped attacks over the exact-label oracle in 2,000 resampled
populations. Rare attacks expose differences hidden by the balanced panel.}
\label{fig:apps-prevalence}
\end{figure}

AUROC and audit loss agree broadly for this frozen mixture, but that need not
hold when score scales or the mix of attacks changes; Appendix~\ref{app:apps-diagnostics}
checks this issue by APPS family and input length. This is an external monitor
test, not a deployment claim: the attacks are stored rather than adapted to each
monitor, all learned representations come from one model, and a stronger code
monitor may beat the tested text baseline.
